\documentclass[10pt,a4paper]{amsart}
\usepackage{amsmath,amssymb,amsthm,graphicx}
\usepackage{verbatim}
\newenvironment{reference}{\comment}{\endcomment}
%\newenvironment{reference}{}{}
\newenvironment{slogan}{\comment}{\endcomment}
\newenvironment{history}{\comment}{\endcomment}
% For commutative diagrams we use Xy-pic
\usepackage[all]{xy}

% We use 2cell for 2-commutative diagrams.
\xyoption{2cell}
% Theorem environments.
%
\theoremstyle{plain}
\newtheorem{theorem}[subsection]{Theorem}
\newtheorem{proposition}[subsection]{Proposition}
\newtheorem{lemma}[subsection]{Lemma}

\theoremstyle{definition}
\newtheorem{definition}[subsection]{Definition}
\newtheorem{example}[subsection]{Example}
\newtheorem{exercise}[subsection]{Exercise}
\newtheorem{situation}[subsection]{Situation}

\theoremstyle{remark}
\newtheorem{remark}[subsection]{Remark}
\newtheorem{remarks}[subsection]{Remarks}

\numberwithin{equation}{subsection}

% Macros
%
\def\lim{\mathop{\mathrm{lim}}\nolimits}
\def\colim{\mathop{\mathrm{colim}}\nolimits}
\def\Spec{\mathop{\mathrm{Spec}}}
\def\Hom{\mathop{\mathrm{Hom}}\nolimits}
\def\Ext{\mathop{\mathrm{Ext}}\nolimits}
\def\SheafHom{\mathop{\mathcal{H}\!\mathit{om}}\nolimits}
\def\SheafExt{\mathop{\mathcal{E}\!\mathit{xt}}\nolimits}
\def\Sch{\mathit{Sch}}
\def\Mor{\mathop{\mathrm{Mor}}\nolimits}
\def\Ob{\mathop{\mathrm{Ob}}\nolimits}
\def\Sh{\mathop{\mathit{Sh}}\nolimits}
\def\NL{\mathop{N\!L}\nolimits}
\def\CH{\mathop{\mathrm{CH}}\nolimits}
\def\proetale{{pro\text{-}\acute{e}tale}}
\def\etale{{\acute{e}tale}}
\def\QCoh{\mathit{QCoh}}
\def\Ker{\mathop{\mathrm{Ker}}}
\def\Im{\mathop{\mathrm{Im}}}
\def\Coker{\mathop{\mathrm{Coker}}}
\def\Coim{\mathop{\mathrm{Coim}}}

% Boxtimes
%
\DeclareMathSymbol{\boxtimes}{\mathbin}{AMSa}{"02}

%
% Macros for moduli stacks/spaces
%
\def\QCohstack{\mathcal{QC}\!\mathit{oh}}
\def\Cohstack{\mathcal{C}\!\mathit{oh}}
\def\Spacesstack{\mathcal{S}\!\mathit{paces}}
\def\Quotfunctor{\mathrm{Quot}}
\def\Hilbfunctor{\mathrm{Hilb}}
\def\Curvesstack{\mathcal{C}\!\mathit{urves}}
\def\Polarizedstack{\mathcal{P}\!\mathit{olarized}}
\def\Complexesstack{\mathcal{C}\!\mathit{omplexes}}
% \Pic is the operator that assigns to X its picard group, usage \Pic(X)
% \Picardstack_{X/B} denotes the Picard stack of X over B
% \Picardfunctor_{X/B} denotes the Picard functor of X over B
\def\Pic{\mathop{\mathrm{Pic}}\nolimits}
\def\Picardstack{\mathcal{P}\!\mathit{ic}}
\def\Picardfunctor{\mathrm{Pic}}
\def\Deformationcategory{\mathcal{D}\!\mathit{ef}}
\usepackage{cleveref}
\crefname{originallemma}{original lemma}{original lemmas}
\Crefname{originallemma}{Original lemma}{Original lemmas}
\makeatletter
\newlabel{lemma-bijective-map}{{lemma-bijective-map}{source}{}{sourcepointer-1}{}}
\newlabel{lemma-bijective-map@cref}{{[originallemma][0][]lemma-bijective-map}{[1][0][]source}{}{}{}}
\makeatother
\crefname{originallemma}{original lemma}{original lemmas}
\Crefname{originallemma}{Original lemma}{Original lemmas}
\makeatletter
\newlabel{lemma-connected-component-intersection-compact-Hausdorff}{{lemma-connected-component-intersection-compact-Hausdorff}{source}{}{sourcepointer-2}{}}
\newlabel{lemma-connected-component-intersection-compact-Hausdorff@cref}{{[originallemma][0][]lemma-connected-component-intersection-compact-Hausdorff}{[1][0][]source}{}{}{}}
\makeatother
\crefname{originallemma}{original lemma}{original lemmas}
\Crefname{originallemma}{Original lemma}{Original lemmas}
\makeatletter
\newlabel{lemma-intersection-closed-in-quasi-compact}{{lemma-intersection-closed-in-quasi-compact}{source}{}{sourcepointer-3}{}}
\newlabel{lemma-intersection-closed-in-quasi-compact@cref}{{[originallemma][0][]lemma-intersection-closed-in-quasi-compact}{[1][0][]source}{}{}{}}
\makeatother
\crefname{originallemma}{original lemma}{original lemmas}
\Crefname{originallemma}{Original lemma}{Original lemmas}
\makeatletter
\newlabel{lemma-inverse-limit-quasi-compact}{{lemma-inverse-limit-quasi-compact}{source}{}{sourcepointer-4}{}}
\newlabel{lemma-inverse-limit-quasi-compact@cref}{{[originallemma][0][]lemma-inverse-limit-quasi-compact}{[1][0][]source}{}{}{}}
\makeatother
\title{Profinite spaces are exactly the Hausdorff quasi-compact totally disconnected spaces, and are cofiltered limits of finite discrete spaces}
\author{The Stacks Project Authors}
\date{stacks-project:topology.tex:local-snapshot-20261012; GNU FDL 1.2 or later}
\begin{document}
\maketitle
\noindent\textit{Editorial note.} All delivered mathematical text is the author's own, copied from the pinned Stacks Project chapter source sources/stacks-topology/pinned-original-source.tex; no statement or proof step was rewritten, reordered or removed.\par\medskip
\noindent\textit{Editorial note.} Preamble provenance: the Stacks Project keeps one single shared preamble for every chapter; this card takes it directly from the staged sources/stacks-topology/preamble.tex as three verbatim spans (lines 8-12, 14-18 and 154-217), so no derived file was written and nothing was added or redefined. The wrapper class and the base packages are supplied by the assembler.\par\medskip
\noindent\textit{Editorial note.} Bibliography provenance: neither the counted statement nor its proof contains a citation (this lane rescans the whole span before writing the card), so the card carries no bibliography and no entry was reconstructed.\par\medskip
\noindent\textit{Editorial note.} Selection provenance: this card was opened under the file-based control contract (CONTROL.md), with the mechanical pre-card gate run on its statement and proof interval before the card was written.\par\medskip
\noindent\textit{Allowed standard prerequisites (editorial).}\begin{itemize}
\item The statement that an inverse limit of quasi-compact spaces is quasi-compact. Invoked inside the counted proof (source lines 4364-4365). Makes the limit of finite discrete spaces quasi-compact.
\item The statement that in a quasi-compact Hausdorff space the connected component of a point is the intersection of its clopen neighbourhoods. Invoked inside the counted proof (source lines 4391-4393). Makes the comparison map injective.
\item The statement that an intersection of closed subsets of a quasi-compact space is nonempty when all finite intersections are. Invoked inside the counted proof (source line 4394). Makes the comparison map surjective.
\item The statement that a bijective map of quasi-compact Hausdorff spaces is a homeomorphism. Invoked at the end of the counted proof (source line 4395). Upgrades the bijection to a homeomorphism.
\end{itemize}\medskip
\ifdefined\hypertarget\hypertarget{sourcepointer-1}{}\fi
\noindent\textit{Source reference.} Original lemma lemma-bijective-map: sources/stacks-topology/pinned-original-source.tex, source lines 3498--3498. Navigation pointer only: the counted text refers to the Stacks item lemma-bijective-map, which is an existing labeled result of this same chapter and is not part of this card. The pointer carries the target's own label as its number and its definition line as its span; its content is not reproduced here.\par\medskip
\ifdefined\hypertarget\hypertarget{sourcepointer-2}{}\fi
\noindent\textit{Source reference.} Original lemma lemma-connected-component-intersection-compact-Hausdorff: sources/stacks-topology/pinned-original-source.tex, source lines 2078--2078. Navigation pointer only: the counted text refers to the Stacks item lemma-connected-component-intersection-compact-Hausdorff, which is an existing labeled result of this same chapter and is not part of this card. The pointer carries the target's own label as its number and its definition line as its span; its content is not reproduced here.\par\medskip
\ifdefined\hypertarget\hypertarget{sourcepointer-3}{}\fi
\noindent\textit{Source reference.} Original lemma lemma-intersection-closed-in-quasi-compact: sources/stacks-topology/pinned-original-source.tex, source lines 1941--1941. Navigation pointer only: the counted text refers to the Stacks item lemma-intersection-closed-in-quasi-compact, which is an existing labeled result of this same chapter and is not part of this card. The pointer carries the target's own label as its number and its definition line as its span; its content is not reproduced here.\par\medskip
\ifdefined\hypertarget\hypertarget{sourcepointer-4}{}\fi
\noindent\textit{Source reference.} Original lemma lemma-inverse-limit-quasi-compact: sources/stacks-topology/pinned-original-source.tex, source lines 2692--2692. Navigation pointer only: the counted text refers to the Stacks item lemma-inverse-limit-quasi-compact, which is an existing labeled result of this same chapter and is not part of this card. The pointer carries the target's own label as its number and its definition line as its span; its content is not reproduced here.\par\medskip
\begin{lemma}
\label{lemma-profinite}
Let $X$ be a topological space.
The following are equivalent
\begin{enumerate}
\item $X$ is a profinite space, and
\item $X$ is Hausdorff, quasi-compact, and totally disconnected.
\end{enumerate}
If this is true, then $X$ is a cofiltered limit of finite discrete
spaces.
\end{lemma}

\begin{proof}
Assume (1). Choose a diagram $i \mapsto X_i$ of finite discrete spaces
such that $X = \lim X_i$. As each $X_i$ is Hausdorff and quasi-compact we find
that $X$ is quasi-compact by Lemma \ref{lemma-inverse-limit-quasi-compact}.
If $x, x' \in X$ are distinct points, then $x$ and $x'$ map to
distinct points in some $X_i$. Hence $x$ and $x'$ have disjoint
open neighbourhoods, i.e., $X$ is Hausdorff. In exactly the same way
we see that $X$ is totally disconnected.

\medskip\noindent
Assume (2). Let $\mathcal{I}$ be the set of finite disjoint union
decompositions $X = \coprod_{i \in I} U_i$ with $U_i$ nonempty open
(and closed) for all $i \in I$.
For each $I \in \mathcal{I}$ there is a continuous map
$X \to I$ sending a point of $U_i$ to $i$. We define a partial
ordering: $I \leq I'$ for $I, I' \in \mathcal{I}$ if and only
if the covering corresponding to $I'$ refines the covering corresponding
to $I$. In this case we obtain a canonical map $I' \to I$. In other
words we obtain an inverse system of finite discrete spaces over $\mathcal{I}$.
The maps $X \to I$ fit together and we obtain a continuous map
$$
X \longrightarrow \lim_{I \in \mathcal{I}} I
$$
We claim this map is a homeomorphism, which finishes the proof.
(The final assertion follows too as the partially ordered set
$\mathcal{I}$ is directed: given two disjoint union decompositions
of $X$ we can find a third refining both.)
Namely, the map is injective as $X$ is totally disconnected
and hence $\{x\}$ is the intersection of all open and closed subsets
of $X$ containing $x$
(Lemma \ref{lemma-connected-component-intersection-compact-Hausdorff})
and  the map is surjective by
Lemma \ref{lemma-intersection-closed-in-quasi-compact}.
By Lemma \ref{lemma-bijective-map} the map is a homeomorphism.
\end{proof}

\end{document}
